\section{Проверка адекватности математических моделей}

\subsection{Понятие адекватности модели}

Математическая модель -- это формализованное описание существенных свойств, связей и закономерностей исследуемого объекта или процесса. При построении модели часть свойств реального объекта неизбежно отбрасывается, поэтому любая математическая модель является \textbf{приближённым} описанием действительности.

\textbf{Адекватностью математической модели} называют её способность с требуемой точностью воспроизводить свойства реального объекта, существенные для поставленной цели исследования, в заданной области применимости модели.

Таким образом, утверждение об адекватности всегда связано с:
\begin{itemize}
	\item целью моделирования;
	\item набором наблюдаемых величин;
	\item допустимым уровнем погрешности;
	\item диапазоном параметров и режимов работы системы.
\end{itemize}

Модель может быть адекватной для одной задачи и неадекватной для другой. Например, упрощённая модель движения тела без сопротивления воздуха может быть достаточной для грубой оценки времени полёта на малых расстояниях, но непригодной для точного расчёта траектории.

Пусть реальный объект характеризуется наблюдаемой величиной
\[
y^{\mathrm{exp}}(x),
\]
а математическая модель даёт прогноз
\[
y^{\mathrm{mod}}(x,\theta),
\]
где $\theta$ --- вектор параметров модели. Тогда величина
\[
r(x,\theta)=y^{\mathrm{exp}}(x)-y^{\mathrm{mod}}(x,\theta)
\]
называется \textbf{остатком}, то есть расхождением модели и наблюдения.

В простейшем случае модель можно считать адекватной, если величина расхождения в рассматриваемой области не превышает заданного допустимого уровня:
\[
\|r\|\leq\varepsilon,
\]
где допустимый уровень $\varepsilon$ должен задаваться исходя из цели моделирования и точности исходных данных.

Однако одного совпадения с отдельным экспериментом недостаточно. Необходимо также проверять корректность реализации модели, устойчивость полученных результатов, воспроизводимость и способность модели описывать данные, которые не использовались при настройке её параметров.

Следует различать несколько близких понятий:
\begin{itemize}
	\item \textbf{верификация} --- правильно ли решены уравнения выбранной модели;
	\item \textbf{валидация} --- правильно ли сама модель описывает реальный объект;
	\item \textbf{калибровка} --- какие значения параметров модели лучше всего согласуют её с наблюдениями.
\end{itemize}

Условно эту разницу можно сформулировать следующим образом:
\[
\boxed{\text{верификация: ``правильно ли мы решаем модель?''}}
\]
\[
\boxed{\text{валидация: ``правильную ли модель мы решаем?''}}
\]

\subsection{Верификация и валидация}

\textbf{Верификация} заключается в проверке правильности математических преобразований, численных алгоритмов и программной реализации модели. При верификации проверяется соответствие вычислительного результата исходной математической постановке.

Основные способы верификации:
\begin{enumerate}
	\item сравнение с аналитическим решением для частного случая;
	\item сравнение с эталонным численным решением;
	\item сравнение результатов различных независимых алгоритмов;
	\item исследование сходимости решения при измельчении сетки или уменьшении шага;
	\item проверка законов сохранения и других инвариантов;
	\item тестирование отдельных частей программного комплекса;
	\item оценка влияния погрешностей округления и критериев остановки итерационных процессов.
\end{enumerate}

Если численный метод имеет порядок сходимости $p$, то в асимптотической области при уменьшении характерного шага сетки $h$ ожидается поведение ошибки вида
\[
e_h\approx Ch^p.
\]
Следовательно,
\[
\frac{e_h}{e_{h/2}}\approx 2^p.
\]
Наблюдение такого закона является одним из способов проверки правильности реализации численного метода.

\textbf{Валидация} --- это проверка соответствия математической модели реальному объекту или процессу. В отличие от верификации, валидация требует данных о \textbf{реальной} системе:
\begin{itemize}
	\item экспериментальных измерений;
	\item натурных наблюдений;
	\item результатов эксплуатации;
	\item достоверных эмпирических зависимостей.
\end{itemize}

Типичная схема проверки модели имеет вид
\[
\text{математическая модель}
\longrightarrow
\text{численный алгоритм}
\longrightarrow
\text{программа}
\longrightarrow
\text{расчёт}
\longrightarrow
\text{сравнение с экспериментом}.
\]

Если расчёт не согласуется с экспериментом, причиной могут быть:
\begin{enumerate}
	\item неправильная математическая модель и/или неверные значения её параметров;
	\item недостаточно точный численный метод;
	\item ошибка программной реализации;
	\item погрешности экспериментальных данных.
\end{enumerate}

Поэтому верификация обычно должна предшествовать валидации: бессмысленно оценивать адекватность физической модели, если неизвестно, правильно ли она решается вычислительным методом.

\subsection{Сопоставление с экспериментальными данными}

Пусть имеется набор экспериментальных наблюдений
\[
\left\{(x_i,y_i^{\mathrm{exp}})\right\}_{i=1}^{N},
\]
а расчёт по модели даёт значения
\[
y_i^{\mathrm{mod}}=y^{\mathrm{mod}}(x_i,\theta).
\]

Вводятся остатки
\[
r_i=y_i^{\mathrm{exp}}-y_i^{\mathrm{mod}}.
\]

Сравнение расчёта и эксперимента включает не только вычисление величины ошибок, но и анализ их структуры.

Хорошая модель должна, как правило, характеризоваться:
\begin{itemize}
	\item малыми остатками;
	\item отсутствием выраженного систематического смещения;
	\item отсутствием закономерной зависимости остатка от входных параметров;
	\item согласованностью величины расхождений с известной погрешностью эксперимента;
	\item правильным воспроизведением качественного поведения системы;
	\item работоспособностью на данных, не использовавшихся при настройке параметров.
\end{itemize}

Важно различать данные, используемые для \textbf{калибровки}, и независимые данные, используемые для \textbf{валидации}. Если параметры модели определяются по набору $D_{\mathrm{cal}}$, то оценивать прогнозирующую способность модели желательно на \textbf{независимом} наборе $D_{\mathrm{val}}$.

Иначе модель может хорошо описывать данные, по которым её параметры были настроены, но плохо предсказывать независимые наблюдения. В статистических моделях такую ситуацию связывают с переобучением.

Для динамических моделей сравниваются не только отдельные значения, но и такие характеристики, как:
\begin{itemize}
	\item положение и высота максимумов;
	\item периоды колебаний;
	\item скорости роста и затухания;
	\item переходные процессы;
	\item стационарные состояния.
\end{itemize}

Сопоставление желательно проводить не в одной точке, а во всей предполагаемой области применимости модели.

\subsection{Критерии качества модели}

Для количественной оценки качества модели применяются различные показатели.

\textbf{Средняя абсолютная ошибка}:
\[
\mathrm{MAE}
=
\frac{1}{N}
\sum_{i=1}^{N}
\left|
y_i^{\mathrm{exp}}-y_i^{\mathrm{mod}}
\right|.
\]

\textbf{Среднеквадратичная ошибка}:
\[
\mathrm{MSE}
=
\frac{1}{N}
\sum_{i=1}^{N}
\left(
y_i^{\mathrm{exp}}-y_i^{\mathrm{mod}}
\right)^2.
\]

\textbf{Корень из среднеквадратичной ошибки}:
\[
\mathrm{RMSE}
=
\sqrt{
	\frac{1}{N}
	\sum_{i=1}^{N}
	\left(
	y_i^{\mathrm{exp}}-y_i^{\mathrm{mod}}
	\right)^2
}.
\]

В отличие от MAE, критерии MSE и RMSE сильнее штрафуют большие ошибки.

Если важна относительная ошибка, используют, например,
\[
\delta_i
=
\frac{
	|y_i^{\mathrm{exp}}-y_i^{\mathrm{mod}}|
}{
	|y_i^{\mathrm{exp}}|
}.
\]
Такой показатель неприменим или становится неустойчивым при $y_i^{\mathrm{exp}}\approx 0$.

Кроме численной величины ошибки необходимо учитывать:
\begin{itemize}
	\item устойчивость результата;
	\item физическую интерпретируемость;
	\item вычислительную стоимость;
	\item сложность модели;
	\item способность модели к прогнозированию.
\end{itemize}

Чрезмерное усложнение модели не всегда улучшает её качество. Если две модели дают сопоставимую точность, обычно предпочтительна более простая модель с меньшим числом параметров.

\subsection{Калибровка параметров}

Во многих моделях имеются параметры, значения которых заранее известны неточно:
\[
\theta=(\theta_1,\ldots,\theta_m).
\]

\textbf{Калибровкой} называют определение параметров модели на основании экспериментальных данных.

Пусть
\[
y_i^{\mathrm{mod}}=y(x_i,\theta).
\]
Тогда параметры можно определять из задачи минимизации
\[
\theta^*
=
\arg\min_{\theta}J(\theta),
\]
где функционал качества, например, имеет вид
\[
J(\theta)
=
\sum_{i=1}^{N}
\left(
y_i^{\mathrm{exp}}-y(x_i,\theta)
\right)^2.
\]

Если измерения имеют различную точность, можно использовать взвешенный функционал:
\[
J(\theta)
=
\sum_{i=1}^{N}
w_i
\left(
y_i^{\mathrm{exp}}-y(x_i,\theta)
\right)^2,
\]
где веса $w_i$ отражают относительную значимость или достоверность отдельных измерений.

В общем случае функционал качества может иметь и другой вид; его выбирают в соответствии с физическим смыслом задачи и требуемым критерием согласования модели с экспериментом.

Оценивать адекватность модели только по тем данным, которые использовались для калибровки, недостаточно. Для валидации желательно использовать независимые данные. Иначе можно получить параметры, хорошо описывающие конкретный набор измерений, но не обеспечивающие требуемую прогнозирующую способность модели.

\subsection{Идентифицируемость модели}

Даже если задача калибровки сформулирована, из экспериментальных данных не всегда можно однозначно определить параметры модели.

\textbf{Идентифицируемость} характеризует возможность восстановления параметров модели по доступным наблюдениям.

Пусть совокупность всех наблюдаемых модельных величин задаётся отображением $y=F(\theta)$.

Если $F(\theta^{(1)})=F(\theta^{(2)})$ влечёт $\theta^{(1)}=\theta^{(2)}$, то отображение параметров в наблюдаемые величины является \textbf{инъективным}, и параметры могут быть однозначно идентифицированы.

Различают:
\begin{itemize}
	\item \textbf{структурную идентифицируемость} --- возможность определить параметры при идеальных, точных данных;
	\item \textbf{практическую идентифицируемость} --- возможность достаточно точно определить параметры по реальным зашумлённым данным.
\end{itemize}

Структурная идентифицируемость может быть \textbf{глобальной}, если набор параметров определяется единственным образом во всей допустимой области, и \textbf{локальной}, если однозначность имеет место лишь в некоторой окрестности рассматриваемого значения параметров.

Модель может быть структурно идентифицируема, но практически плохо идентифицируема, если разные значения параметров приводят к почти одинаковым результатам.

Для локального анализа вводится матрица чувствительности
\[
S_{ij}
=
\frac{\partial y_i}{\partial\theta_j}.
\]

Если её столбцы линейно зависимы или почти линейно зависимы, то воздействие некоторых параметров на результат трудно различить. В таком случае задача определения параметров становится плохо обусловленной.

Проблемы практической идентифицируемости могут быть ослаблены:
\begin{itemize}
	\item проведением дополнительных экспериментов;
	\item измерением дополнительных величин;
	\item фиксацией части параметров;
	\item изменением постановки эксперимента;
	\item упрощением модели;
	\item применением регуляризации.
\end{itemize}

При этом регуляризация не устраняет структурную неидентифицируемость, а лишь стабилизирует задачу оценивания параметров.

\subsection{Анализ чувствительности}

\textbf{Анализ чувствительности} исследует изменение результата модели при изменении её входных данных и параметров.

Если выход модели
\[
y=F(\theta_1,\ldots,\theta_m),
\]
то локальная чувствительность к параметру $\theta_j$ определяется производной
\[
S_j
=
\frac{\partial F}{\partial\theta_j}.
\]

Для сравнения параметров различных размерностей удобно использовать безразмерный коэффициент чувствительности:
\[
S_j^{\mathrm{rel}}
=
\frac{\theta_j}{F}
\frac{\partial F}{\partial\theta_j}.
\]

Он показывает относительное изменение результата при малом относительном изменении параметра.

На практике производная может оцениваться конечной разностью:
\[
\frac{\partial F}{\partial\theta_j}
\approx
\frac{
	F(\theta+\Delta\theta_j e_j)-F(\theta)
}{
	\Delta\theta_j
}.
\]

Различают:
\begin{itemize}
	\item \textbf{локальный анализ чувствительности}, исследующий малые изменения около некоторого номинального набора параметров;
	\item \textbf{глобальный анализ чувствительности}, исследующий влияние параметров во всём допустимом диапазоне. Обычно требует варьирования нескольких параметров одновременно во всей допустимой области.
\end{itemize}

Анализ чувствительности позволяет:
\begin{enumerate}
	\item определить наиболее важные параметры модели;
	\item выявить практически несущественные параметры;
	\item понять причины большой неопределённости результата;
	\item спланировать необходимые измерения;
	\item обнаружить проблемы идентифицируемости;
	\item упростить модель.
\end{enumerate}

\subsection{Учёт неопределённости}

Результат математического моделирования обычно содержит неопределённость, связанную как с самой моделью, так и с входными данными.

Основные источники неопределённости:
\begin{enumerate}
	\item неопределённость параметров;
	\item погрешности измерений;
	\item неопределённость начальных и граничных условий;
	\item упрощения математической модели;
	\item численная погрешность;
	\item случайный характер моделируемого процесса.
\end{enumerate}

Особое значение имеет \textbf{структурная погрешность модели}: даже при точно известных параметрах и практически безошибочном численном решении сама выбранная математическая модель может систематически отличаться от реального объекта вследствие принятых упрощений и неучтённых эффектов.

Часто различают два основных типа неопределённости.

\textbf{Алеаторная неопределённость} (объективная случайность) обусловлена объективной случайностью процесса, например случайными внешними воздействиями. Увеличение числа измерений позволяет точнее оценить её статистические характеристики, но не устраняет саму случайность процесса.

\textbf{Эпистемическая неопределённость} (недостаток знаний) обусловлена недостатком знаний: неточностью параметров, неполнотой модели или недостаточным объёмом экспериментальных данных. Она потенциально может быть уменьшена дополнительными исследованиями.

\subsection{Резюме}

Проверка адекватности математической модели представляет собой не один критерий, а совокупность взаимосвязанных процедур:
\[
\boxed{
	\begin{array}{c}
		\text{верификация (правильно решаем модель?)}
		\\
		\downarrow
		\\
		\text{калибровка параметров (какие параметры согласуют модель с экспериментом?)}
		\\
		\downarrow
		\\
		\text{валидация по независимым данным (правильно ли модель описывает реальность?)}
		\\
		\downarrow
		\\
		\text{анализ чувствительности и неопределённости}
		\\
		\downarrow
		\\
		\text{определение области применимости модели}
	\end{array}
}
\]

На практике эти этапы могут выполняться итерационно: неудовлетворительные результаты валидации приводят к уточнению модели, её параметров или численного алгоритма.

Модель следует считать адекватной не вообще, а только относительно заданной цели, набора наблюдаемых характеристик, требуемой точности и установленной области применимости.

\newpage